\chapter{Operações Lógicas II}

\begin{questao}{Questão 13}

Sejam $p$ e $q$ duas proposições. Se $V(\sim p$ $\underline{\vee}$ $q) = V$, logo podemos afirmar que

\begin{alternativas}
\item $V(p) = F$
\item $V(p) = V$
\item $p$ e $q$ têm valorações iguais.
\item Não é possível afirmar nada sobre a valoração das proposições.
\end{alternativas}

\end{questao}

\begin{comentario_questao}

\comentarioitem{A}{Para que $V(\sim p$ $\underline{\vee}$ $q) = V$, é necessário que exatamente uma das proposições $\sim p$ ou $q$ seja verdadeira. Isso não garante que $V(p) = F$, pois também pode ocorrer o caso em que $V(\sim p) = F$ e $V(q) = V$, ou seja, $V(p)$ pode ser $V$.}

\comentarioitem{B}{Como visto, não é possível garantir que $V(p) = V$, pois existem casos em que $p$ pode ser falso e ainda assim a disjunção exclusiva ser verdadeira.}

\comentarioitem{C}{Essa é alternativa correta. Para que $V(\sim p$ $\underline{\vee}$ $q) = V$, é necessário que exatamente uma das proposições $\sim p$ ou $q$ seja verdadeira. Isso implica que $\sim p$ e $q$ têm valorações opostas, ou seja, se $p$ é verdadeiro, então $q$ é verdadeiro também; e quando $p$ é falso, $q$ é falso. Portanto, $p$ e $q$ têm valorações iguais.}

\comentarioitem{D}{Sim, é possível afirmar que $p$ e $q$ têm valorações iguais.}

\respostacorreta{C}

\end{comentario_questao}

\newpage

\begin{questao}{Questão 14}

Sejam as proposições $p$: Está frio'' e $q$: Está chovendo''. Em linguagem natural, \\$\sim p$ $\rightarrow \sim q$ é traduzida como

\begin{alternativas}
\item Se não está calor então não está chovendo.
\item Se não está frio então não está chovendo.
\item É falso que se não está calor então não está chovendo.
\item Não é verdade que se está frio então não está chovendo.
\end{alternativas}

\end{questao}

\begin{comentario_questao}
    \comentarioitem{A}{``Está calor'' não é necessariamente a negação de ``Está frio''. É possível que esteja nublado, por exemplo, o que não se encaixa em nenhuma das proposições. Portanto, essa alternativa não é correta.}

    \comentarioitem{B}{Essa é a alternativa correta. A proposição $\sim p$ $\rightarrow \sim q$ é lida como ``Se não está frio, então não está chovendo''.}

    \comentarioitem{C}{``Está calor'' não é necessariamente a negação de ``Está frio''. E a expressão ``É falso que'' negaria a expressão inteira, algo como $\sim (\ldots)$. Portanto, essa alternativa não é correta.}

    \comentarioitem{D}{Essa é a negação da proposição inteira. Em notação simbólica seria $\sim (p$ $\rightarrow \sim q)$. Portanto, essa alternativa não é correta.}

    \respostacorreta{B}
\end{comentario_questao}

\newpage

\begin{questao}{Questão 15}

Sejam $p$, $q$ e $r$ proposições quaisquer. Se $V(p \wedge q) = F$, logo podemos afirmar {\bf categoricamente} que

\begin{alternativas}
\item $V((p \wedge q) \rightarrow r) = V$
\item $V(p \vee q) = F$
\item $V(\sim p) = V$
\item $V(r \rightarrow q) = F$
\end{alternativas}

\end{questao}

\begin{comentario_questao}

    \comentarioitem{A}{Essa é a alternativa correta. A proposição $p \wedge q$ é falsa, o que significa que $V(F \rightarrow r) = V$ para qualquer valor de $r$.}

    \comentarioitem{B}{A proposição $p \vee q$ pode ser verdadeira ou falsa, dependendo dos valores de $p$ e $q$. Se ambos forem falsos, então $p \vee q$ será falso; se pelo menos um deles for verdadeiro, então $p \vee q$ será verdadeiro. Portanto, não podemos afirmar categoricamente que $V(p \vee q) = F$.}

    \comentarioitem{C}{A proposição $\sim p$ pode ser verdadeira ou falsa, dependendo do valor de $p$. Se $p$ for verdadeiro, então $\sim p$ será falso; se $p$ for falso, então $\sim p$ será verdadeiro. Portanto, não podemos afirmar categoricamente que $V(\sim p) = V$.}

    \comentarioitem{D}{A proposição $r \rightarrow q$ pode ser verdadeira ou falsa, dependendo dos valores de $r$ e $q$. Se $r$ for verdadeiro e $q$ for falso, então $r \rightarrow q$ será falso; caso contrário, será verdadeiro. Portanto, não podemos afirmar categoricamente que $V(r \rightarrow q) = F$.}

    \respostacorreta{A}
\end{comentario_questao}

\newpage

\begin{questao}{Questão 16}

Sejam as proposições $p$: ``Jorge é rico'' e $q$: ``Carlos é feliz''. Em linguagem simbólica,

``Jorge não é rico, se e somente se, Carlos for feliz''

é expressa como

\begin{alternativas}
\item $\sim p \leftrightarrow q$
\item $\sim p \leftrightarrow \sim q$
\item $\sim p \leftarrow q$
\item $\sim p \leftarrow \sim q$
\end{alternativas}

\end{questao}

\begin{comentario_questao}

    \comentarioitem{A}{Essa é a alternativa correta. A proposição ``Jorge não é rico, se e somente se, Carlos for feliz'' é expressa como $\sim p \leftrightarrow q$.}

    \comentarioitem{B}{A proposição $\sim p \leftrightarrow \sim q$ seria lida como ``Jorge não é rico, se e somente se, Carlos não for feliz'', o que não corresponde à frase dada.}

    \comentarioitem{C}{A proposição $\sim p \leftarrow q$ seria lida como ``Se Carlos for feliz, então Jorge não é rico'', o que não corresponde à frase dada.}

    \comentarioitem{D}{A proposição $\sim p \leftarrow \sim q$ seria lida como ``Se Carlos não for feliz, então Jorge não é rico'', o que não corresponde à frase dada.}

    \respostacorreta{A}
\end{comentario_questao}

\newpage

\begin{questao}{Questão 17}

Sejam $p$, $q$ e $r$ proposições quaisquer. Se $V(p \wedge q) = F$, logo podemos afirmar {\bf categoricamente} que

\begin{alternativas}
\item $V((p \wedge q) \leftrightarrow r) = V$
\item $V(p \rightarrow q) = F$
\item $V(r \leftrightarrow r) = V$
\item $V(r \rightarrow (p \wedge q)) = F$
\end{alternativas}

\end{questao}

\begin{comentario_questao}

    \comentarioitem{A}{O valor de $(p \wedge q) \leftrightarrow r$ pode ser falso, no caso de $r$ ser verdadeiro, ou seja, $V((p \wedge q) \leftrightarrow r) = F$. Portanto, não podemos afirmar categoricamente que $V((p \wedge q) \leftrightarrow r) = V$.}

    \comentarioitem{B}{A proposição $p \rightarrow q$ pode ser verdadeira ou falsa, dependendo dos valores de $p$ e $q$. Se $p$ for verdadeiro e $q$ for falso, então $p \rightarrow q$ será falso; caso contrário, será verdadeiro. Portanto, não podemos afirmar categoricamente que $V(p \rightarrow q) = F$.}

    \comentarioitem{C}{A proposição $r \leftrightarrow r$ é uma tautologia, ou seja, é sempre verdadeira independentemente do valor de $r$. Portanto, essa alternativa é correta.}

    \comentarioitem{D}{A proposição $r \rightarrow (p \wedge q)$ pode ser verdadeira ou falsa, dependendo dos valores de $r$, $p$ e $q$. 
    \begin{itemize}
        \item Se $r$ for verdadeiro, então $r \rightarrow (p \wedge q)$ será falso (pois $p \wedge q$ é falso);
        \item Se $r$ for falso, então $r \rightarrow (p \wedge q)$ pode ser tanto verdadeiro quanto falso.
    \end{itemize} }

    \respostacorreta{C}
\end{comentario_questao}

\newpage

\begin{questao}{Questão 18}

Sejam $p$, $q$ e $r$ proposições quaisquer. Se $V(p$ $\underline{\vee}$ $q) = V$ e $V(r) = F$, logo podemos afirmar {\bf categoricamente} que

\begin{alternativas}
\item $V(p \leftrightarrow q) = V$
\item $V\big((p \leftrightarrow q)$ $\underline{\vee}$ $r\big) = V$
\item $V\big((p$ $\underline{\vee}$ $q) \leftrightarrow r\big) = F$
\item $V\big(r \rightarrow (p$ $\underline{\vee}$ $q)\big) = F$
\end{alternativas}

\end{questao}

\begin{comentario_questao}

    \comentarioitem{A}{Se $V(p$ $\underline{\vee}$ $q) = V$, isso significa que exatamente uma das proposições $p$ ou $q$ é verdadeira. Portanto, $p$ e $q$ têm valorações opostas, o que implica que $V(p \leftrightarrow q) = F$. Portanto, não podemos afirmar categoricamente que $V(p \leftrightarrow q) = V$.}

    \comentarioitem{B}{ A proposição $(p \leftrightarrow q)$ é falsa, como visto na alternativa A. E a disjunção exclusiva de uma proposição falsa com uma proposição falsa ($r$) é falsa. Portanto, $V\big((p \leftrightarrow q)$ $\underline{\vee}$ $r\big) = F$. Portanto, não podemos afirmar categoricamente que $V\big((p \leftrightarrow q)$ $\underline{\vee}$ $r\big) = V$.}   

    \comentarioitem{C}{A proposição $(p$ $\underline{\vee}$ $q) \leftrightarrow r$ é falsa, pois $V(p$ $\underline{\vee}$ $q) = V$ e $V(r) = F$. Portanto, podemos afirmar categoricamente que $V\big((p$ $\underline{\vee}$ $q) \leftrightarrow r\big) = F$.}

    \comentarioitem{D}{A proposição $r \rightarrow (p$ $\underline{\vee}$ $q)$ é verdadeira, pois $V(r) = F$. Portanto, $V\big(r \rightarrow (p$ $\underline{\vee}$ $q)\big) = V$. Portanto, está incorreta a afirmação de que $V\big(r \rightarrow (p$ $\underline{\vee}$ $q)\big) = F$.}

    \respostacorreta{C}
\end{comentario_questao}    